\documentclass[11pt]{article}
\usepackage[margin=1in]{geometry}
\usepackage{amsmath,amssymb,amsthm,enumitem}
\usepackage{hyperref}

\newtheorem{theorem}{Theorem}[section]
\newtheorem{lemma}[theorem]{Lemma}
\newtheorem{corollary}[theorem]{Corollary}
\newtheorem{proposition}[theorem]{Proposition}
\newtheorem{definition}[theorem]{Definition}
\newtheorem{assumption}[theorem]{Assumption}
\newtheorem{remark}[theorem]{Remark}
\newtheorem{example}[theorem]{Example}

\newcommand{\adj}{\operatorname{adj}}
\newcommand{\spann}{\operatorname{span}}
\newcommand{\Ran}{\operatorname{Ran}}
\newcommand{\Vresp}{\mathcal V_{\mathrm{resp}}}
\renewcommand{\S}{\mathfrak S}
\newcommand{\E}{\mathfrak E}

\title{Section 0: Operational Foundations of the Sector Cut\\
\large A prefix to \emph{Three Central Theorems of the
Sector-Resolved Response Classification}}
\author{}
\date{}

\begin{document}
\maketitle

\begin{abstract}
This document supplies the material that the companion theorems
(\texttt{three\_theorems\_proofs.tex}, Theorems I--III) presuppose but do
not themselves supply: a physical definition of what counts as a
\emph{sector} and what counts as a \emph{sector cut}. We define a sector
operationally, as the equivalence class of admissible selective
dissipative interventions sharing a common retained Riesz projector, and
we define the ideal sector cut as the strong-intervention
($\kappa\to\infty$) limit of such an intervention. Theorem 0A below is
obtained from the existing protected-channel theorem (Theorem III of the
companion document) by a role substitution
$(\Gamma,D,B)\mapsto(\kappa,D_S,A_\Gamma)$; no new estimate is required.
Theorems 0B, 0C, 0E and Proposition 0D are short consequences of 0A. The
only genuinely new technical content is the sharpened statement of kernel
universality (Theorem 0B and Lemma 0B$'$), which identifies exactly when
agreement of $\ker D_S$ alone (rather than of the full oblique projector
$P_S$) suffices for two cut implementations to agree.
\end{abstract}

% ==================================================================
\section{Operational experiment specification}
% ==================================================================

\begin{definition}[Experimental specification]
\label{def:E}
An \emph{experimental specification} is a tuple
\[
\E=(\mathcal L_0,\rho_0,\mathcal V_p,O,\Gamma,K)
\]
where $\mathcal L_0$ is an unperturbed Lindbladian on a finite-dimensional
Liouville space $\mathcal V\cong\mathbb C^N$, $\rho_0$ is a reference
steady state ($\mathcal L_0\rho_0=0$), $\mathcal V_p$ is a probe-injection
superoperator, $O$ is a readout observable, $\Gamma\ge0$ is the native
dissipation scale entering $\mathcal L_0$ as in
\eqref{eq:AGamma-recall} below, and $K\subset\Omega$ is a compact
observation window in the frequency domain $\Omega$ on which the weak-probe
transfer function is analytic.

The associated linear response is
\[
\chi_\Gamma(z)=\langle\!\langle O\mid A_\Gamma(z)^{-1}\mid c\rangle\!\rangle,
\qquad
A_\Gamma(z)=\Gamma D+B(z),
\qquad
\lvert c\rangle\!\rangle=\mathcal V_p\lvert\rho_0\rangle\!\rangle,
\label{eq:AGamma-recall}
\tag{0.1}
\]
matching the generator family \eqref{eq:AGamma} of the companion document
with $p^\dagger=\langle\!\langle O\mid$.
\end{definition}

\begin{definition}[Admissible cut generator]
\label{def:admissible}
Let $\E$ be an experimental specification and let $S$ be a designated
target sector. A superoperator $\mathcal C_S$ (equivalently, its linear
response representation $D_S$ acting on $\mathcal V$) is an
\emph{admissible cut generator} for $S$ if:
\begin{enumerate}[label=(C\arabic*),nosep]
\item \textbf{GKSL admissibility.} $\mathcal L_0+\kappa\mathcal C_S$ is a
      valid Lindblad generator (completely positive, trace-preserving
      semigroup generator) for every $\kappa\ge0$.
\item \textbf{Selective damping.} The nonzero spectrum of $D_S$ damps the
      target sector and leaves the subspace to be retained inside
      $\ker D_S$.
\item \textbf{Semisimple zero mode.} $0$ is a semisimple eigenvalue of
      $D_S$, so that the Riesz projection
      \[
      P_S=\frac1{2\pi i}\oint_{\gamma}(\zeta I-D_S)^{-1}\,d\zeta,
      \qquad Q_S=I-P_S,
      \]
      with $\gamma$ a small contour around $0$ enclosing no other
      eigenvalue of $D_S$, is well defined (Lemma~\ref{lem:riesz-recall}
      below, restated from the companion document).
\item \textbf{Noninvasive reference condition.} $\mathcal C_S(\rho_0)=0$.
\item \textbf{Fixed preparation and readout.} $\mathcal V_p$ and $O$ are
      unchanged by the cut.
\end{enumerate}
\end{definition}

\begin{lemma}[Riesz splitting, restated]
\label{lem:riesz-recall}
Under (C3), $P_S^2=P_S$, $P_SD_S=D_SP_S=0$,
$\mathcal V=P_S\mathcal V\oplus Q_S\mathcal V$ (a direct, generally
non-orthogonal sum), $\Ran P_S=\ker D_S$, and $D_S|_{Q_S\mathcal V}$ is
invertible.
\end{lemma}

\begin{proof}
Identical to Lemma~\ref{lem:riesz} of the companion document with $D$
replaced by $D_S$; the argument uses only the holomorphic functional
calculus applied to $D_S$ and does not reference the dissipation scale
$\Gamma$.
\end{proof}

\begin{definition}[Operational sector]
\label{def:sector}
Two admissible cut generators $\mathcal C_S,\mathcal C_S'$ (for the same
target $S$, possibly with different microscopic implementations) are
\emph{equivalent}, written $\mathcal C_S\sim\mathcal C_S'$, iff their
Riesz projections agree:
\[
\mathcal C_S\sim\mathcal C_S'
\quad:\Longleftrightarrow\quad
P_S=P_S'.
\]
The \emph{operational sector} is the equivalence class
$\S=[\mathcal C_S]$. We call $Q_S=I-P_S$ the \emph{suppressed part} and
$P_S\mathcal V$ the \emph{retained part} of $\S$.
\end{definition}

\begin{remark}
Equivalence is stated at the level of the full (generally oblique) Riesz
projector $P_S$, not merely of the kernel $\ker D_S$. Two admissible
generators with the same kernel but different range decompositions
($\Ran D_S\ne\Ran D_S'$) give different oblique projectors and hence, in
general, different cut limits; see Lemma~\ref{lem:kernel-suffices} for
the (physically important) case in which the kernel alone does suffice.
\end{remark}

% ==================================================================
\section{Theorem 0A: the operational cut limit}
% ==================================================================

Fix $\E$ and $\Gamma$, and let $\mathcal C_S$ be admissible with Riesz
data $(P_S,Q_S)$ as above. Write $D_S$ for its response-space
representation and form the two-scale family
\[
A_{\Gamma,\kappa}(z)=A_\Gamma(z)+\kappa D_S,
\qquad \kappa\ge0,\ z\in\Omega,
\]
with sign convention such that the nonzero eigenvalues of $D_S$ have
positive real part (so that $\kappa D_S$ damps $Q_S\mathcal V$ as
$\kappa\to\infty$), consistent with \S4.1 of the revision note.

\begin{theorem}[0A: Operational cut limit]
\label{thm:0A}
Assume, on the compact window $K$:
\begin{enumerate}[label=(H0\arabic*),nosep]
\item $0$ is a semisimple eigenvalue of $D_S$ (Definition~\ref{def:admissible}, (C3));
\item $Q_SD_SQ_S$ is invertible on $Q_S\mathcal V$ (automatic from
      Lemma~\ref{lem:riesz-recall});
\item the \emph{retained block}
      $A_{P_S,\Gamma}(z):=P_SA_\Gamma(z)P_S|_{P_S\mathcal V}$ is invertible
      for every $z\in K$, with $\sup_{z\in K}\|A_{P_S,\Gamma}(z)^{-1}\|<\infty$;
\item $\Gamma$ is fixed and $X_{K,\Gamma}:=\sup_{z\in K}\|A_\Gamma(z)\|<\infty$
      (automatic from analyticity and boundedness on compacts,
      Assumption~1.2 of the companion document).
\end{enumerate}
Then there is $\kappa_*(\Gamma)<\infty$ such that for $\kappa>\kappa_*(\Gamma)$
and $z\in K$,
\[
\chi_{\Gamma,\kappa}(z):=p^\dagger A_{\Gamma,\kappa}(z)^{-1}c
=\chi_{\mathrm{cut},\Gamma}^{(S)}(z)+\kappa^{-1}G_1(z)+O(\kappa^{-2})
\]
uniformly on $K$, where
\[
\boxed{\;
\chi_{\mathrm{cut},\Gamma}^{(S)}(z)
=p^\dagger P_S\bigl[P_SA_\Gamma(z)P_S\bigr]^{-1}P_Sc
\;}
\label{eq:0A}
\tag{0.2}
\]
and $G_1$ is given by the block formula \eqref{eq:F1} of the companion
document with the substitution $B\mapsto A_\Gamma$, $D_Q\mapsto Q_SD_SQ_S$,
$P\mapsto P_S$, $Q\mapsto Q_S$. Moreover $\kappa_*(\Gamma)=O(\Gamma)$ as
$\Gamma\to\infty$, since $X_{K,\Gamma}=O(\Gamma)$.
\end{theorem}

\begin{proof}
This is Theorem III of the companion document applied verbatim under the
role substitution $(\Gamma,D,B(z))\mapsto(\kappa,D_S,A_\Gamma(z))$ with
$\Gamma$ held fixed as a parameter of the (now $\kappa$-independent)
coherent part $A_\Gamma(z)$. Hypotheses (H01)--(H03) above are exactly
(H1)--(H3) of Theorem~\ref{thm:III} under this substitution; (H04) is
(H4) of Theorem~\ref{thm:III}, and it is automatic here because $\Gamma$
is fixed on the window $K$ and $B$, $D$ are individually bounded on
compacts by Assumption~1.2. The block computation of
Theorem~\ref{thm:III} (Riesz splitting $\mathcal V=P_S\mathcal V\oplus
Q_S\mathcal V$, Schur complement of the retained block, Neumann series in
$\kappa^{-1}$ for the $(\Gamma D_S+\text{const})^{-1}$-type block) is
reproduced without modification, yielding \eqref{eq:0A} as the order-zero
term and the stated $G_1$ as the order-$\kappa^{-1}$ term. Finally,
$\kappa_*(\Gamma)$ plays the role of $\Gamma_*$ in the proof of
Theorem~\ref{thm:III}, where it was shown to be controlled by
$\|D_Q^{-1}\|\sup_K\|B\|\,\|Q\|^2$; under the substitution this becomes
$\|(Q_SD_SQ_S)^{-1}\|\,X_{K,\Gamma}\,\|Q_S\|^2$, and $X_{K,\Gamma}\le
\Gamma\|D\|+\sup_K\|B(\cdot)\|=O(\Gamma)$ by \eqref{eq:AGamma-recall}.
\end{proof}

\begin{remark}[The two limits do not commute for free]
\label{rem:double-limit}
Because $\kappa_*(\Gamma)=O(\Gamma)$, uniformity of the $\kappa$-expansion
on $K$ degrades as $\Gamma\to\infty$: Theorem~\ref{thm:0A} guarantees the
expansion only once $\kappa/\Gamma$ is large enough (a constant multiple
of $\|(Q_SD_SQ_S)^{-1}\|\,\|D\|\,\|Q_S\|^2$, up to the bounded part of
$B$). Consequently
\[
\lim_{\Gamma\to\infty}\lim_{\kappa\to\infty}\chi_{\Gamma,\kappa}(z)
\quad\text{and}\quad
\lim_{\kappa\to\infty}\lim_{\Gamma\to\infty}\chi_{\Gamma,\kappa}(z)
\]
are not shown equal by this theorem; the standing convention of this
theory (Section~\ref{sec:convention}) fixes the left-hand order as the
definition of the classified quantity, and Gate U7 is reserved for probing
the (possibly nonempty) region where the two orders disagree.
\end{remark}

% ==================================================================
\section{Theorem 0B: kernel universality}
% ==================================================================

\begin{theorem}[0B: Kernel universality]
\label{thm:0B}
Let $\mathcal C_S,\mathcal C_S'$ be admissible cut generators for the same
target $S$ satisfying $P_S=P_S'$. Then, under the hypotheses of
Theorem~\ref{thm:0A} applied to each,
\[
\chi_{\mathrm{cut},\Gamma}^{[\mathcal C_S]}(z)
=\chi_{\mathrm{cut},\Gamma}^{[\mathcal C_S']}(z)
\qquad\text{for all }z\in K.
\]
That is, the ideal sector-cut response depends on the admissible cut
generator only through the equivalence class $\S=[\mathcal C_S]$ of
Definition~\ref{def:sector}, not on the dephasing/damping rate, jump
operator normalization, or internal basis of the suppressed sector.
\end{theorem}

\begin{proof}
By \eqref{eq:0A}, $\chi_{\mathrm{cut},\Gamma}^{(S)}(z)$ is a function of
$(P_S,A_\Gamma(z),c,p)$ alone; it does not otherwise reference $D_S$.
Since $P_S=P_S'$ by hypothesis and $A_\Gamma(z),c,p$ are the same for both
cuts, the right-hand sides of \eqref{eq:0A} for $\mathcal C_S$ and
$\mathcal C_S'$ coincide termwise.
\end{proof}

\begin{lemma}[Kernel alone suffices for self-adjoint cuts]
\label{lem:kernel-suffices}
Suppose $D_S$ is self-adjoint with respect to the Hilbert--Schmidt inner
product on operator space, $\langle X,D_SY\rangle_{\mathrm{HS}}
=\langle D_SX,Y\rangle_{\mathrm{HS}}$ for all $X,Y$ (as holds, e.g., for
$D_S$ built from $\mathcal D[L]$ with a Hermitian jump operator $L=L^\dagger$,
i.e.\ pure dephasing). Then:
\begin{enumerate}[label=(\roman*),nosep]
\item $0$ is automatically a semisimple eigenvalue of $D_S$ (self-adjoint
      operators on a finite-dimensional inner-product space are
      diagonalizable with orthogonal eigenspaces), so (C3) is automatic;
\item the Riesz projection $P_S$ is the \emph{orthogonal} projection onto
      $\ker D_S$, and is therefore determined by $\ker D_S$ alone: if
      $D_S,D_S'$ are both self-adjoint admissible cut generators with
      $\ker D_S=\ker D_S'$, then $P_S=P_S'$, and Theorem~\ref{thm:0B}
      applies.
\end{enumerate}
\end{lemma}

\begin{proof}
(i) The spectral theorem for finite-dimensional self-adjoint operators
gives an orthogonal eigenbasis and real eigenvalues; the eigenspace for
$0$ equals $\ker D_S$ with no generalized (nilpotent) part, i.e.\ the zero
eigenvalue is semisimple.
(ii) For self-adjoint $D_S$, $\Ran D_S=(\ker D_S)^{\perp_{\mathrm{HS}}}$
(orthogonal complement of the kernel with respect to the Hilbert--Schmidt
inner product), so the direct-sum decomposition
$\mathcal V=\ker D_S\oplus\Ran D_S$ of
Lemma~\ref{lem:riesz-recall} is orthogonal and the associated Riesz
projection $P_S$ coincides with the orthogonal projector onto
$\ker D_S$, which depends on $\ker D_S$ alone. Hence
$\ker D_S=\ker D_S'\Rightarrow P_S=P_S'$.
\end{proof}

\begin{remark}
Lemma~\ref{lem:kernel-suffices} is the precise sense in which the
informal phrase ``kernel-defined operational uniqueness'' of the revision
note is correct: it holds within the self-adjoint (e.g.\ pure-dephasing)
class of admissible cut generators, but not for general non-self-adjoint
$D_S$, where the oblique range $\Ran D_S$ must also be matched. The
minimal $\Lambda$-model verification of Section~11 of the revision note
uses a Hermitian dephasing jump operator and therefore falls inside this
safe class.
\end{remark}

% ==================================================================
\section{Theorem 0C: covariant representation invariance}
% ==================================================================

\begin{theorem}[0C: Covariant representation invariance]
\label{thm:0C}
Let $T$ be invertible and transform the full experimental data covariantly:
\[
(A_\Gamma,D_S,c,p,P_S)
\;\longmapsto\;
(TA_\Gamma T^{-1},\,TD_ST^{-1},\,Tc,\,T^{-\dagger}p,\,TP_ST^{-1}).
\]
Then $\chi_{\mathrm{full},\Gamma}$, $\chi_{\mathrm{cut},\Gamma}^{(S)}$,
$R_{S,\Gamma}=\chi_{\mathrm{full},\Gamma}-\chi_{\mathrm{cut},\Gamma}^{(S)}$,
and the index $\nu[R_S]$ of Definition~\ref{def:nu-recall} are all
unchanged.
\end{theorem}

\begin{proof}
$\chi_{\mathrm{full},\Gamma}$ is invariant by Lemma~1(iii) of the
companion document (representation invariance of $F_\Gamma$). For
$\chi_{\mathrm{cut},\Gamma}^{(S)}$, first note that the Riesz projection
transforms covariantly under similarity:
\[
TP_ST^{-1}
=T\left(\frac1{2\pi i}\oint_\gamma(\zeta I-D_S)^{-1}d\zeta\right)T^{-1}
=\frac1{2\pi i}\oint_\gamma\bigl(\zeta I-TD_ST^{-1}\bigr)^{-1}d\zeta,
\]
using $T(\zeta I-D_S)^{-1}T^{-1}=(\zeta I-TD_ST^{-1})^{-1}$; hence the
transformed projector is exactly the Riesz projector of the transformed
generator $TD_ST^{-1}$, confirming that $P_S\mapsto TP_ST^{-1}$ is the
\emph{correct} (not merely notational) transformation law. Substituting
into \eqref{eq:0A},
\[
(T^{-\dagger}p)^\dagger(TP_ST^{-1})\bigl[(TP_ST^{-1})(TA_\Gamma T^{-1})(TP_ST^{-1})\bigr]^{-1}(TP_ST^{-1})(Tc)
\]
collapses by the same telescoping cancellation as in Lemma~1(i) of the
companion document (each adjacent $T^{-1}T$ pair cancels) to
$p^\dagger P_S[P_SA_\Gamma P_S]^{-1}P_Sc$, i.e.\ $\chi_{\mathrm{cut},\Gamma}^{(S)}$
is unchanged. Invariance of $R_{S,\Gamma}$ follows by subtraction, and
invariance of $\nu[R_S]$ follows because $\nu$ (Definition~7.1 of the
master response package) is defined purely in terms of the invariant
function $R_{S,\Gamma}$.
\end{proof}

\begin{remark}
The point of stating $P_S\mapsto TP_ST^{-1}$ explicitly is that a
transformation which fixes the sector \emph{label} but mixes sector and
complement coordinates (i.e.\ does not conjugate $P_S$ along with
everything else) does \emph{not} fall under Theorem~\ref{thm:0C}: it
represents a different physical intervention, not a change of coordinates
of the same intervention. Gate U4 (arbitrary-coordinate negative control)
verifies that such non-covariant mixing does change $R_S$, as it must.
\end{remark}

% ==================================================================
\section{Proposition 0D: frozen/self-consistent decomposition}
% ==================================================================

\begin{proposition}[0D: Frozen/self-consistent decomposition]
\label{prop:0D}
Let $\chi_{\mathrm{cut}}^{\mathrm{fr}}$ denote the sector-cut response
computed with $\rho_0$, $c$, $p$ frozen at their full-model values, and
let $\chi_{\mathrm{cut}}^{\mathrm{sc}}$ denote the response computed after
re-solving for the steady state of the cut generator
$\mathcal L_0+\kappa\mathcal C_S$ (in the $\kappa\to\infty$ limit) and
updating $c=\mathcal V_p\rho_0^{\mathrm{cut}}$ accordingly. Then
\[
\chi_{\mathrm{full}}-\chi_{\mathrm{cut}}^{\mathrm{sc}}
=\underbrace{\chi_{\mathrm{full}}-\chi_{\mathrm{cut}}^{\mathrm{fr}}}_{R_S^{\mathrm{path}}}
+\underbrace{\chi_{\mathrm{cut}}^{\mathrm{fr}}-\chi_{\mathrm{cut}}^{\mathrm{sc}}}_{R_S^{\mathrm{reprep}}}.
\]
\end{proposition}

\begin{proof}
Add and subtract $\chi_{\mathrm{cut}}^{\mathrm{fr}}$; this is an algebraic
identity requiring no hypothesis.
\end{proof}

\begin{lemma}[Noninvasive condition forces agreement]
\label{lem:0Dprime}
If $\mathcal C_S(\rho_0)=0$ (Definition~\ref{def:admissible}, (C4)) and,
for every $\kappa\ge0$, $\mathcal L_0+\kappa\mathcal C_S$ has a unique
steady state, then $\rho_0^{\mathrm{cut}}=\rho_0$ for every $\kappa\ge0$
and, in the $\kappa\to\infty$ limit, $R_S^{\mathrm{reprep}}=0$, i.e.\
$\chi_{\mathrm{cut}}^{\mathrm{fr}}=\chi_{\mathrm{cut}}^{\mathrm{sc}}$.
\end{lemma}

\begin{proof}
$\mathcal L_0(\rho_0)=0$ by definition of $\rho_0$ as the reference
steady state, and $\mathcal C_S(\rho_0)=0$ by (C4); hence, by linearity,
$(\mathcal L_0+\kappa\mathcal C_S)(\rho_0)=0$ for every $\kappa\ge0$, so
$\rho_0$ is \emph{a} steady state of the cut generator at every $\kappa$.
By the assumed uniqueness of the steady state at each $\kappa$, it is
\emph{the} steady state, i.e.\ $\rho_0^{\mathrm{cut}}(\kappa)=\rho_0$ for
all $\kappa\ge0$; the $\kappa\to\infty$ limit of a constant sequence is
that constant, so $\rho_0^{\mathrm{cut}}=\rho_0$ persists in the ideal cut
limit. Consequently $c=\mathcal V_p\rho_0^{\mathrm{cut}}=\mathcal V_p\rho_0$
is the same source used to compute $\chi_{\mathrm{cut}}^{\mathrm{fr}}$, so
the two constructions coincide and $R_S^{\mathrm{reprep}}=0$ in
Proposition~\ref{prop:0D}.
\end{proof}

% ==================================================================
\section{Theorem 0E: class robustness under cut ambiguity}
% ==================================================================

\begin{theorem}[0E: Class robustness]
\label{thm:0E}
Let $\chi_{\mathrm{cut}}^{(1)},\chi_{\mathrm{cut}}^{(2)}$ be two finite-strength
(not necessarily ideal, $\kappa<\infty$) implementations of a sector cut
for the same operational sector $\S$, satisfying, uniformly on $K$,
\[
\bigl\|\chi_{\mathrm{cut}}^{(1)}-\chi_{\mathrm{cut}}^{(2)}\bigr\|_K=O(\Gamma^{-q})
\]
as $\Gamma\to\infty$. Suppose the true sector-mediated response for
implementation (2) obeys, uniformly on $K$,
\[
R_S^{(2)}(z)=\Gamma^{-\nu}r_\nu(z)+o(\Gamma^{-\nu}),
\qquad r_\nu\not\equiv0,
\]
with $q>\nu$. Then $R_S^{(1)}$ has the same leading order and the same
leading coefficient:
\[
R_S^{(1)}(z)=\Gamma^{-\nu}r_\nu(z)+o(\Gamma^{-\nu}).
\]
In particular, the finite decision algorithm (Theorem~\ref{thm:algorithm}
of the companion document) returns the same class $\nu$ from either
implementation.
\end{theorem}

\begin{proof}
$R_S^{(1)}-R_S^{(2)}=\chi_{\mathrm{cut}}^{(2)}-\chi_{\mathrm{cut}}^{(1)}
=O(\Gamma^{-q})$ uniformly on $K$ by hypothesis (the full response
$\chi_{\mathrm{full}}$ is common to both and cancels). Adding this to the
assumed expansion of $R_S^{(2)}$,
\[
R_S^{(1)}=R_S^{(2)}+O(\Gamma^{-q})
=\Gamma^{-\nu}r_\nu+o(\Gamma^{-\nu})+O(\Gamma^{-q}).
\]
Since $q>\nu$, $O(\Gamma^{-q})=o(\Gamma^{-\nu})$, and it absorbs into the
error term without altering the leading coefficient.
\end{proof}

% ==================================================================
\section{Standing convention on the two dissipation scales}
\label{sec:convention}
% ==================================================================

\begin{definition}[Order of limits]
\label{def:convention}
Throughout the theory, the sector-cut response at fixed $\Gamma$ is
\emph{defined} as
\[
\chi_{\mathrm{cut},\Gamma}(z):=\lim_{\kappa\to\infty}\chi_{\Gamma,\kappa}(z)
\]
(Theorem~\ref{thm:0A}), and the native-dissipation classification of
$R_{S,\Gamma}=\chi_{\mathrm{full},\Gamma}-\chi_{\mathrm{cut},\Gamma}^{(S)}$
into $\nu\in\{\infty\}\cup(0,\infty)\cup\{0\}$ (Theorems I--III of the
companion document) is then carried out by letting $\Gamma\to\infty$.
That is, the standing convention is
\[
\lim_{\Gamma\to\infty}\Bigl[\lim_{\kappa\to\infty}\chi_{\Gamma,\kappa}(z)\Bigr],
\]
\emph{not} the reverse order. Remark~\ref{rem:double-limit} shows the
reverse order is not established by Theorem~\ref{thm:0A} and is left as
an open question (Gate U7).
\end{definition}

% ==================================================================
\section*{Dependency chart}
% ==================================================================

\begin{center}
\begin{tabular}{lll}
Lemma~\ref{lem:riesz-recall} (companion Lemma~2, restated)
 & $\longrightarrow$ & Theorem~\ref{thm:0A}\\[2pt]
Companion Theorem~III (protected channel)
 & $\longrightarrow$ & Theorem~\ref{thm:0A} (role substitution)\\[2pt]
Theorem~\ref{thm:0A}
 & $\longrightarrow$ & Theorem~\ref{thm:0B}, Remark~\ref{rem:double-limit}\\[2pt]
Spectral theorem for self-adjoint operators
 & $\longrightarrow$ & Lemma~\ref{lem:kernel-suffices}\\[2pt]
Companion Lemma~1 (basis invariance) + Riesz covariance
 & $\longrightarrow$ & Theorem~\ref{thm:0C}\\[2pt]
Linearity + steady-state uniqueness
 & $\longrightarrow$ & Lemma~\ref{lem:0Dprime} $\longrightarrow$ Proposition~\ref{prop:0D}\\[2pt]
Asymptotic comparison ($q>\nu$)
 & $\longrightarrow$ & Theorem~\ref{thm:0E}\\[2pt]
Theorems~\ref{thm:0A}--\ref{thm:0E}
 & $\longrightarrow$ & companion Theorems I--III apply to
 $\chi_{\mathrm{cut},\Gamma}^{(S)}$ of \eqref{eq:0A} as the operational
 realization of the sector cut
\end{tabular}
\end{center}

\end{document}
